% Related work, drafted with an assistant and not yet edited.
% The fixture for the paper branch: the markup section does NOT apply here.

\section{Related Work}
\label{sec:related}

Reinforcement learning has emerged as a comprehensive framework for adaptive
traffic signal control, and a substantial body of work underscores its potential
in this domain~\cite{genders2016,wei2019}. Early value-based methods serve as
the foundation for much of what followed, providing a robust basis for the
actor-critic designs that came later~\cite{mnih2015}.

Tabular $Q$-learning remains the canonical starting point, offering an
interpretability that deep methods sacrifice, and it is crucial for establishing
baselines on small intersection topologies~\cite{abdulhai2003}. Function
approximation was subsequently introduced, thereby enabling generalisation
across unseen states and significantly improving sample efficiency. Notably, the
resulting methods exhibited stable convergence on networks of 4--9
intersections, highlighting the importance of careful reward shaping.

Unlike prior studies that optimise average delay in isolation, this work
emphasises the interplay between queue length and phase stability---a
distinction that has received comparatively little attention.

Despite its comprehensive coverage of the single-intersection case, the
literature faces several challenges when scaled to corridors. Observers have
noted that coordination overhead may potentially dominate at higher flow
densities, although the effect is robust to sensor dropout in the two studies
that measured it~\cite{chu2020}. Coordination across corridors therefore remains
an open problem. Our two-branch design provides an efficient and reliable
framework for addressing it.
